% Source: Lee, "Introduction to Riemannian Manifolds" (2nd ed.), Chapter 4 "Connections", PDF pages 98-126 of the cited source.
% Transcribed from the cited source, with manual checks for a few pages, and organized according to leanblueprint
% conventions (semantic \label, bracketed titles, \uses dependency edges, \dcref to Lee's own numbering) below.

\section{Differentiating Vector Fields}

For a smooth curve $\gamma:I\to\mathbb R^n$,
\begin{equation}
\gamma'(t)=\dot\gamma^i(t)\frac{\partial}{\partial x^i}\bigg|_{\gamma(t)}.
\label{eq:4.1}
\end{equation}
\begin{equation}
\gamma''(t)=\ddot\gamma^i(t)\frac{\partial}{\partial x^i}\bigg|_{\gamma(t)}.
\label{eq:4.2}
\end{equation}
Dots denote derivatives in $t$. A curve is a straight line exactly when it admits a parametrization with $\gamma''=0$.

For $Y=Y^i\partial_i$ on $\mathbb R^n$ and $v=v^i\partial_i|_p$, the Euclidean directional derivative is
\[
\bar\nabla_vY=v(Y^i)\partial_i|_p,
\qquad v(Y^i)=v^j\frac{\partial Y^i(p)}{\partial x^j}.
\]
For a vector field $X$, this gives
\begin{equation}
\bar\nabla_XY=X(Y^i)\frac{\partial}{\partial x^i}.\tag{4.3}
\end{equation}

If $M\subset\mathbb R^n$ is an embedded submanifold, the tangential acceleration of a curve is $\gamma''{}^\top=\pi^\top(\gamma'')$. For a local field $Y$ on $M$, extend it to $\widetilde Y$ near $M$ and define
\begin{equation}
\nabla_v^\top Y=\pi^\top(\bar\nabla_v\widetilde Y).\tag{4.4}
\end{equation}
The projection is onto $T_pM$; independence of the extension is part of
\cref{prop:embedded-submanifold-covariant-derivative}.
\uses{prop:normal-bundle}

On an abstract manifold, velocity still has the coordinate-independent expression \eqref{eq:4.1}, but $\gamma'(t+h)$ and $\gamma'(t)$ lie in different tangent spaces, so their difference quotient is undefined. Coordinate expressions can therefore give incompatible accelerations (the unit circle has Cartesian acceleration $-\cos t\,\partial_x-\sin t\,\partial_y$ but polar expression $r''\partial_r+\theta''\partial_\theta=0$). A connection supplies a coordinate-independent rule for comparing and differentiating such vectors.

\section{Connections}

\begin{definition}[Connection]
\label{def:connection}
\lean{LeeLib.Ch04.Connection, LeeLib.Ch04.covariantDeriv}
\leanok


Let $\pi:E\to M$ be a smooth vector bundle over a smooth manifold with or without boundary. A connection is a map
\[
\nabla:\mathfrak X(M)\times\Gamma(E)\to\Gamma(E)
\]
such that
\begin{enumerate}
\item[(i)] $\nabla_{f_1X_1+f_2X_2}Y=f_1\nabla_{X_1}Y+f_2\nabla_{X_2}Y$;
\item[(ii)] $\nabla_X(a_1Y_1+a_2Y_2)=a_1\nabla_XY_1+a_2\nabla_XY_2$;
\item[(iii)] $\nabla_X(fY)=f\nabla_XY+(Xf)Y$.
\end{enumerate}
Here $f,f_i\in C^\infty(M)$, $a_i\in\mathbb R$, and the fields have the indicated domains.
\end{definition}

\begin{lemma}[Locality]
\label{lem:connection-is-local-operator}
\lean{LeeLib.Ch04.covariantDeriv_local}
\uses{def:connection}
\leanok
\dcref{ch4:4.1}



If $\nabla$ is a connection, then $\nabla_XY|_p$ depends only on the germs of $X,Y$ near $p$: if $X=\widetilde X$ and $Y=\widetilde Y$ near $p$, then $\nabla_XY|_p=\nabla_{\widetilde X}\widetilde Y|_p$.
\end{lemma}
\begin{proof}
\leanok
\sketch If $Y$ vanishes near $p$, choose a bump function $\varphi$ supported there with $\varphi(p)=1$. Since $\varphi Y=0$, the product rule gives
\begin{equation}
0=\nabla_X(\varphi Y)=(X\varphi)Y+\varphi\nabla_XY.\tag{4.5}
\end{equation}
Both terms except $\nabla_XY|_p$ vanish at $p$. The argument for changing $X$ is analogous, using $C^\infty$-linearity in the first slot.
\end{proof}

\begin{proposition}[Restriction of a Connection]\label{prop:connection-restriction}\dcref{ch4:4.3}
\uses{lem:connection-is-local-operator}

If $\nabla$ is a connection on $E\to M$ and $U\subseteq M$ is open, there is a unique connection $\nabla^U$ on $E|_U$ satisfying
\begin{equation}
\nabla^U_{X|_U}(Y|_U)=(\nabla_XY)|_U.\tag{4.6}
\end{equation}
\end{proposition}
\begin{proof}
\sketch Extend local fields and sections near each $p\in U$ to global ones with bump functions and define
\begin{equation}
\nabla^U_XY|_p=(\nabla_{\widetilde X}\widetilde Y)|_p.\tag{4.7}
\end{equation}
Lemma \ref{lem:connection-is-local-operator} makes this independent of extensions and smooth; the connection axioms follow locally.
\end{proof}

\begin{proposition}\label{prop:covariant-deriv-dependence-at-point}\dcref{ch4:4.5}
\lean{LeeLib.Ch04.covariantDeriv_dependence_at_point}
\uses{lem:connection-is-local-operator}
\leanok



For a connection, $\nabla_XY|_p$ depends only on the germ of $Y$ near $p$ and the single value $X_p$.
\end{proposition}
\begin{proof}
\leanok
\sketch Locality handles $Y$. If $X_p=0$, restrict to coordinates and write $X=X^i\partial_i$ with $X^i(p)=0$; restriction to the coordinate neighborhood and $C^\infty$-linearity give $\nabla_XY|_p=X^i(p)\nabla_{\partial_i}Y|_p=0$.
\end{proof}

Thus for $v\in T_pM$ and a local section $Y$, define $\nabla_vY$ by extending $v$ to a local field $X$ and setting $\nabla_vY=\nabla_XY|_p$; the preceding proposition makes this well defined.

\section{Connections in the Tangent Bundle}

A connection in $TM$ is a connection on $M$ and is a map $\mathfrak X(M)^2\to\mathfrak X(M)$ satisfying (i)--(iii). It is not a $(1,2)$-tensor because the second slot obeys the product rule rather than $C^\infty(M)$-linearity.

For a local frame $(E_i)$, define connection coefficients by
\begin{equation}
\nabla_{E_i}E_j=\Gamma^k_{ij}E_k.\tag{4.8}
\end{equation}
\begin{proposition}
\label{prop:connection-coefficients-frame}
\lean{LeeLib.Ch04.connectionCoeff, LeeLib.Ch04.covariantDeriv_frame_eq_sum_connectionCoeff, LeeLib.Ch04.covariantDeriv_eq_connectionCoeff_formula}
\uses{def:connection}
\leanok
\dcref{ch4:4.6}



If $X=X^iE_i$ and $Y=Y^jE_j$ on a local frame domain, then
\begin{equation}
\nabla_XY=\left(X(Y^k)+X^iY^j\Gamma^k_{ij}\right)E_k.\tag{4.9}
\end{equation}
\end{proposition}
\begin{proof}
\sketch Expand $\nabla_X(Y^jE_j)$ using the product rule and $C^\infty$-linearity in $X$, substitute (4.8), and rename the dummy index in the derivative term.
\end{proof}

\begin{proposition}[Transformation Law for Connection Coefficients]
\label{prop:transformation-law-connection-coefficients}
\lean{LeeLib.Ch04.connectionCoeff_transformation_law}
\uses{def:connection, prop:connection-coefficients-frame}
\leanok
\dcref{ch4:4.7}



If $\widetilde E_i=A^j_iE_j$ are two local frames, their coefficients satisfy
\[
\widetilde\Gamma^k_{ij}=(A^{-1})^k_pA^q_iA^r_j\Gamma^p_{qr}+(A^{-1})^k_pA^q_iE_q(A^p_j).
\]
\end{proposition}
\begin{proof}
\leanok
\sketch Expand $\nabla_{\widetilde E_i}(A^r_jE_r)$ using the product rule from Definition \ref{def:connection}; substitute (4.8) and Proposition \ref{prop:connection-coefficients-frame}, then write $E_p=(A^{-1})^k_p\widetilde E_k$ and $\widetilde E_i(A^p_j)=A^q_iE_q(A^p_j)$. Reusing Proposition \ref{prop:connection-coefficients-frame} gives the stated formula; the derivative term shows that coefficients are not tensor components.
\end{proof}

\section{Existence of Connections}

\begin{example}[The Euclidean Connection]
\label{ex:euclidean-connection}
\dcref{ch4:4.8}
On $T\mathbb R^n$, the formula (4.3) defines the Euclidean connection $\bar\nabla$; its standard-frame coefficients are zero.
\end{example}

\begin{example}[The Tangential Connection on a Submanifold of $\mathbb{R}^n$]
\label{ex:tangential-connection-submanifold}
\uses{prop:normal-bundle, ex:euclidean-connection, prop:extension-smooth-function, prop:embedded-submanifold-covariant-derivative, prop:adapted-orthonormal-frames, exer:extension-smooth-function, prob:embedded-submanifold-covariant-derivative}
\dcref{ch4:4.9}

For an embedded $M\subseteq\mathbb R^n$, define
\[
\nabla_X^TY=\pi^T\bigl(\widetilde\nabla_{\widetilde X}\widetilde Y|_M\bigr),
\]
where extensions are taken near $M$, $\widetilde\nabla$ is the Euclidean connection, and $\pi^T$ is the tangential projection. This is independent of extensions by (4.4) and is a connection.
\end{example}
\begin{proof}
\sketch Linearity in $X,Y$ is immediate. For an extension $\widetilde f$ of $f$,
\[
\nabla_X^T(fY)=\pi^T\bigl(\widetilde\nabla_{\widetilde X}(\widetilde f\widetilde Y)|_M\bigr)=(Xf)Y+f\nabla_X^TY.
\]
\end{proof}

\begin{lemma}
\label{lem:connections-from-connection-coefficients}
\lean{LeeLib.Ch04.connOfCoeff, LeeLib.Ch04.connectionCoeff_connOfCoeff, LeeLib.Ch04.covariantDeriv_eq_connectionCoeff_formula}
\uses{prop:connection-coefficients-frame}
\leanok
\dcref{ch4:4.10}


If a smooth $n$-manifold with or without boundary has a global frame $(E_i)$, then (4.9) gives a one-to-one correspondence between connections on $TM$ and $n^3$ smooth functions $\Gamma^k_{ij}$.
\end{lemma}
\begin{proof}
\sketch A connection supplies $\Gamma^k_{ij}$ through (4.8). Conversely, (4.9) defines a smooth operation satisfying linearity; the product rule is the direct expansion of $X(Y^j)$.
\end{proof}

\begin{proposition}
\label{prop:tangent-bundle-admits-connection}
\uses{lem:connections-from-connection-coefficients}
\dcref{ch4:4.12}

Every smooth manifold with or without boundary admits a connection on its tangent bundle.
\end{proposition}
\begin{proof}
\sketch Choose local connections $\nabla^\alpha$ on a coordinate cover and a subordinate partition of unity $(\varphi_\alpha)$. Set
\begin{equation}
\nabla_XY=\sum_\alpha\varphi_\alpha\nabla_X^\alpha Y.\tag{4.11}
\end{equation}
Local finiteness gives smoothness. The partition identity $\sum_\alpha\varphi_\alpha=1$ yields
\[
\nabla_X(fY)=\sum_\alpha\varphi_\alpha((Xf)Y+f\nabla_X^\alpha Y)=(Xf)Y+f\nabla_XY.
\]
\end{proof}

\begin{proposition}[The Difference Tensor]
\label{prop:difference-tensor}
\lean{LeeLib.Ch04.differenceTensor, LeeLib.Ch04.differenceTensor_apply}
\uses{def:connection}
\leanok
\dcref{ch4:4.13}



For connections $\nabla^0,\nabla^1$ on $TM$, define $D(X,Y)=\nabla_X^1Y-\nabla_X^0Y$. Then $D$ is $C^\infty(M)$-bilinear and hence a $(1,2)$-tensor field.
\end{proposition}
\begin{proof}
\leanok
\sketch First-slot linearity is inherited. Expanding $D(X,fY)$ with the two product rules cancels the two $(Xf)Y$ terms, leaving $D(X,fY)=fD(X,Y)$.
\end{proof}

\begin{theorem}
\label{thm:connections-affine-space}
\lean{LeeLib.Ch04.addOneForm, LeeLib.Ch04.eq_addOneForm_differenceTensor}
\uses{prop:difference-tensor}
\leanok
\dcref{ch4:4.14}



For any connection $\nabla^0$ on a smooth manifold, the set of all connections is the affine space
\[
\mathcal A(TM)=\{\nabla^0+D:D\in\Gamma(T^{(1,2)}TM)\},
\]
where $(\nabla^0+D)_XY=\nabla^0_XY+D(X,Y)$.
\end{theorem}
\begin{proof}
\leanok
\uses{prop:difference-tensor}
\sketch Proposition \ref{prop:difference-tensor} writes any $\nabla$ as $\nabla^0+D$. Conversely, adding a tensor $D$ preserves the two linearities, and $C^\infty$-linearity of $D$ in $Y$ leaves the product rule of $\nabla^0$ unchanged.
\end{proof}
\section{Covariant Derivatives of Tensor Fields}

\begin{proposition}
\label{prop:connection-on-tensor-bundles}
\lean{LeeLib.Ch04.covDerivCovector, LeeLib.Ch04.covDerivCovector_apply, LeeLib.Ch04.covDeriv_pairing_leibniz}
\dcref{ch4:4.15}

Let $M$ be a smooth manifold with or without boundary and let $\nabla$ be a
connection in $TM$. There is a unique induced connection on every tensor
bundle $T^{(k,l)}TM$, denoted again by $\nabla$, satisfying:
\begin{enumerate}
\item[(i)] on $TM$, it is the given connection;
\item[(ii)] on $T^{(0,0)}TM=M\times\mathbb R$, $\nabla_Xf=Xf$;
\item[(iii)] $\nabla_X(F\otimes G)=(\nabla_XF)\otimes G+F\otimes(\nabla_XG)$;
\item[(iv)] $\nabla_X(\operatorname{tr}F)=\operatorname{tr}(\nabla_XF)$ for every
contraction of one covariant and one contravariant index.
\end{enumerate}
It also satisfies
\begin{enumerate}
\item[(a)]
\[
\nabla_X\langle\omega,Y\rangle
=\langle\nabla_X\omega,Y\rangle+\langle\omega,\nabla_XY\rangle;
\]
\item[(b)] for $F\in\Gamma(T^{(k,l)}TM)$, smooth $1$-forms
$\omega^1,\ldots,\omega^k$, and smooth vector fields $Y_1,\ldots,Y_l$,
\begin{align}
(\nabla_XF)(\omega^1,\ldots,\omega^k,Y_1,\ldots,Y_l)
&=X\big(F(\omega^1,\ldots,\omega^k,Y_1,\ldots,Y_l)\big)\nonumber\\
&\quad-\sum_{i=1}^kF(\omega^1,\ldots,\nabla_X\omega^i,\ldots,\omega^k,Y_1,\ldots,Y_l)\nonumber\\
&\quad-\sum_{j=1}^lF(\omega^1,\ldots,\omega^k,Y_1,\ldots,\nabla_XY_j,\ldots,Y_l).
\label{eq:4.12}\tag{4.12}
\end{align}
\end{enumerate}
\end{proposition}
\begin{proof}
\sketch Pairing is the contraction of $\omega\otimes Y$, so (i)--(iv) give
(a). Repeatedly expressing the evaluation of $F$ as contractions of
$F\otimes\omega^1\otimes\cdots\otimes\omega^k\otimes Y_1\otimes\cdots\otimes
Y_l$ gives (b). Formula (a) with functions determines the derivative of every
1-form:
\begin{equation}
(\nabla_X\omega)(Y)=X(\omega(Y))-\omega(\nabla_XY).
\label{eq:4.13}\tag{4.13}
\end{equation}
Then (4.12) determines all higher tensor derivatives, proving uniqueness.
Conversely, define the derivative of $1$-forms by (4.13) and all other
derivatives by (4.12). The terms obtained by differentiating a coefficient
$f$ cancel when an argument is replaced by $f\omega^i$ or $fY_j$, so the
result is tensorial and smooth. Linearity and the tensor-product rule follow
from the corresponding rules for differentiation of functions.
\end{proof}

\begin{proposition}
\label{prop:covariant-derivative-components}
\dcref{ch4:4.16}
Let $M$ be a smooth manifold with or without boundary, let $\nabla$ be a
connection in $TM$, and let $(E_i)$ be a local frame with dual coframe
$(\epsilon^j)$ and coefficients $\nabla_{E_i}E_j=\Gamma^k_{ij}E_k$.
If $X=X^jE_j$, then for
$\omega=\omega_i\epsilon^i$,
\[
\nabla_X\omega=\bigl(X(\omega_k)-X^j\omega_i\Gamma^i_{jk}\bigr)\epsilon^k.
\]
For a tensor
\[
F=F^{i_1\cdots i_k}_{j_1\cdots j_l}
E_{i_1}\otimes\cdots\otimes E_{i_k}\otimes
\epsilon^{j_1}\otimes\cdots\otimes\epsilon^{j_l},
\]
\[
\begin{aligned}
\nabla_XF={}&\Bigl(X(F^{i_1\cdots i_k}_{j_1\cdots j_l})
+\sum_{s=1}^kX^mF^{i_1\cdots p\cdots i_k}_{j_1\cdots j_l}\Gamma^{i_s}_{mp}\\
&\qquad-\sum_{s=1}^lX^mF^{i_1\cdots i_k}_{j_1\cdots p\cdots j_l}\Gamma^p_{ms}\Bigr)
E_{i_1}\otimes\cdots\otimes E_{i_k}\otimes
\epsilon^{j_1}\otimes\cdots\otimes\epsilon^{j_l}.
\end{aligned}
\]
\end{proposition}

\begin{proposition}[The Total Covariant Derivative]
\label{prop:total-covariant-derivative}
\dcref{ch4:4.17}
Let $M$ be a smooth manifold with or without boundary and $\nabla$ a
connection. For $F\in\Gamma(T^{(k,l)}TM)$ define
\[
\nabla F\colon\underbrace{\Omega^1(M)\times\cdots\times\Omega^1(M)}_{k}
\times\underbrace{\mathfrak{X}(M)\times\cdots\times\mathfrak{X}(M)}_{l+1}
\longrightarrow C^\infty(M)
\]
by
\begin{equation}
(\nabla F)(\omega^1,\ldots,\omega^k,Y_1,\ldots,Y_l,X)
=(\nabla_XF)(\omega^1,\ldots,\omega^k,Y_1,\ldots,Y_l).
\tag{4.14}
\end{equation}
This is a smooth $(k,l+1)$-tensor field, called the total covariant derivative.
\end{proposition}
\begin{proof}
\sketch The tensor characterization lemma applies because $\nabla_XF$ is
tensorial in its $k+l$ arguments and a connection is $C^\infty(M)$-linear in
$X$.
\end{proof}

Writing the differentiation index last, if $Y=Y^iE_i$ and
$\omega=\omega_i\epsilon^i$, then
\[
\nabla Y=Y^i_{;j}E_i\otimes\epsilon^j,
\qquad Y^i_{;j}=E_jY^i+Y^k\Gamma^i_{jk},
\]
and
\[
\nabla\omega=\omega_{i;j}\epsilon^i\otimes\epsilon^j,
\qquad \omega_{i;j}=E_j\omega_i-\omega_k\Gamma^k_{ji}.
\]

\begin{proposition}
\label{prop:total-covariant-derivative-components}
\dcref{ch4:4.18}
Let $M$ be a smooth manifold with or without boundary, let $\nabla$ be a
connection in $TM$, and let $(E_i)$ be a smooth local frame with connection
coefficients $\Gamma^k_{ij}$. For a $(k,l)$-tensor $F$,
\[
\begin{aligned}
F^{i_1\cdots i_k}_{j_1\cdots j_l;m}
={}&E_m\bigl(F^{i_1\cdots i_k}_{j_1\cdots j_l}\bigr)
+\sum_{s=1}^kF^{i_1\cdots p\cdots i_k}_{j_1\cdots j_l}\Gamma^{i_s}_{mp}\\
&-\sum_{s=1}^lF^{i_1\cdots i_k}_{j_1\cdots p\cdots j_l}\Gamma^p_{mj_s}.
\end{aligned}
\]
\end{proposition}

\section{Second Covariant Derivatives}

For a smooth manifold $M$ with or without boundary and a connection $\nabla$,
set $\nabla^2F=\nabla(\nabla F)$ for $F\in\Gamma(T^{(k,l)}TM)$ and define
\[
\nabla^2_{X,Y}F(\ldots)=\nabla^2F(\ldots,Y,X).
\]
The reversal of $X,Y$ records that the differentiation index is appended last;
$\nabla^2_{X,Y}$ differentiates first in the $Y$ direction and then in the
$X$ direction. It is $C^\infty(M)$-linear in $Y$, unlike
$\nabla_X(\nabla_YF)$.

\begin{proposition}
\label{prop:second-covariant-derivative-formula}
\lean{LeeLib.Ch04.secondCovariantDeriv_function}
\uses{prop:connection-on-tensor-bundles}
\dcref{ch4:4.21}


For every smooth vector or tensor field $F$,
\[
\nabla^2_{X,Y}F=\nabla_X(\nabla_YF)-\nabla_{\nabla_XY}F.
\]
\end{proposition}
\begin{proof}
\sketch Contract $\nabla F\otimes Y$ in its final two slots to obtain
$\nabla_YF$, and contract $\nabla^2F\otimes X$ and then $Y$ to obtain
$\nabla^2_{X,Y}F$. Applying $\nabla_X$, using compatibility with contraction
and the tensor-product rule, gives
\[
\nabla_X(\nabla_YF)=\nabla^2_{X,Y}F+\nabla_{\nabla_XY}F.
\]
Rearrange to obtain the formula.
\end{proof}

\begin{example}[The Covariant Hessian]
\label{ex:covariant-hessian}
\lean{LeeLib.Ch04.covariantHessianTensor, LeeLib.Ch04.covariantHessianTensor_apply, LeeLib.Ch04.covariantHessian_apply}
\uses{prop:second-covariant-derivative-formula}
\dcref{ch4:4.22}



For $u\in C^\infty(M)$, $\nabla u=du$ as a $1$-form, since
$\nabla_Xu=Xu=du(X)$. The tensor $\nabla^2u=\nabla(du)$ is the covariant
Hessian, and
\[
\nabla^2u(Y,X)=\nabla^2_{X,Y}u
=\nabla_X(\nabla_Yu)-\nabla_{\nabla_XY}u
=X(Yu)-(\nabla_XY)u.
\]
In local coordinates,
\[
\nabla^2u=u_{ij}\,dx^i\otimes dx^j,
\qquad u_{ij}=\partial_j\partial_i u-\Gamma^k_{ji}\partial_ku.
\]
\end{example}

\section{Vector and Tensor Fields Along Curves}

Let $M$ be a smooth manifold with or without boundary and let
$\gamma\colon I\to M$ be smooth. A vector field along $\gamma$ is a
continuous map $V\colon I\to TM$ with $V(t)\in T_{\gamma(t)}M$; it is smooth
when smooth as a map to $TM$. Write $\mathfrak{X}(\gamma)$ for the smooth
fields. Pointwise operations make it a real vector space and a
$C^\infty(I)$-module:
\[
(fX)(t)=f(t)X(t).
\]
The velocity $\gamma'$ is a smooth field along $\gamma$. In $\mathbb R^2$, if
$R$ is rotation by $\pi/2$, then $N=R\gamma'$ has coordinates
$N(t)=(-\dot\gamma^2(t),\dot\gamma^1(t))$.

If $\overline V$ is a smooth vector field on a neighborhood of $\gamma(I)$,
then $V(t)=\overline V_{\gamma(t)}$ is smooth and is called \emph{extendible}.
If $\gamma(t_1)=\gamma(t_2)$ but $\gamma'(t_1)\ne\gamma'(t_2)$, then $\gamma'$
is not extendible. A tensor field along $\gamma$ is defined analogously in
$T^{(k,l)}TM$, and extendibility means restriction of a smooth tensor field
defined near $\gamma(I)$.

\begin{example}
\label{ex:figure-eight-curve}
\dcref{ch4:4.23}
The curve
\[
\gamma\colon(-\pi,\pi)\to\mathbb R^2,
\qquad \gamma(t)=(\sin 2t,\sin t),
\]
is an injective smooth immersion, but its velocity field is not extendible.
\end{example}

\section{Covariant Derivatives Along Curves}

\begin{theorem}[Covariant Derivative Along a Curve]
\label{thm:covariant-derivative-along-curve}
\lean{LeeLib.Ch04.covariantDerivAlongChart, LeeLib.Ch04.covariantDerivAlongChart_add, LeeLib.Ch04.covariantDerivAlongChart_smul}
\uses{ex:figure-eight-curve}
\leanok
\dcref{ch4:4.24}


For a connection $\nabla$ on a smooth manifold $M$ with or without boundary and
a smooth curve $\gamma\colon I\to M$, there is a unique operator
\[
D_t\colon\mathfrak{X}(\gamma)\to\mathfrak{X}(\gamma)
\]
such that
\begin{enumerate}
\item[(i)] $D_t(aV+bW)=aD_tV+bD_tW$ for $a,b\in\mathbb R$;
\item[(ii)] $D_t(fV)=f'V+fD_tV$ for $f\in C^\infty(I)$;
\item[(iii)] if $V$ is the restriction of a smooth field $\widetilde V$ near
$\gamma(I)$, then
\[
D_tV(t)=\nabla_{\gamma'(t)}\widetilde V.
\]
\end{enumerate}
The same construction applies to smooth tensor fields of every type along
$\gamma$.
\end{theorem}
\begin{proof}
\sketch At $t_0\in I$, locality reduces the construction to a coordinate
neighborhood of $\gamma(t_0)$. Writing
$V(t)=V^i(t)\partial_i|_{\gamma(t)}$ gives the forced formula
\begin{equation}
D_tV(t)=\bigl(\dot V^k(t)+\dot\gamma^i(t)V^j(t)\Gamma^k_{ij}(\gamma(t))\bigr)
\partial_k|_{\gamma(t)}.\tag{4.15}
\end{equation}
This proves uniqueness. Define $D_t$ by (4.15) in a chart containing the curve;
the coordinate formula satisfies (i)--(iii), and uniqueness makes definitions
agree on overlaps. At an endpoint, extend the coordinate curve to a slightly
larger interval before restricting.
\end{proof}

\begin{proposition}
\label{prop:covariant-derivative-depends-on-curve}
\uses{prop:covariant-deriv-dependence-at-point, thm:covariant-derivative-along-curve}
\dcref{ch4:4.26}

Let $M$ be a smooth manifold with or without boundary, let $p\in M$,
$v\in T_pM$, and let $Y,\overline Y$ be smooth vector fields
that agree along a smooth curve $\gamma$ with $\gamma(t_0)=p$ and
$\gamma'(t_0)=v$. Then
\[
\nabla_vY=\nabla_v\overline Y.
\]
\end{proposition}
\begin{proof}
\sketch The common restriction $Z(t)=Y_{\gamma(t)}=\overline Y_{\gamma(t)}$
is a field along $\gamma$. Both $Y$ and $\overline Y$ extend $Z$, so property
(iii) of Theorem \ref{thm:covariant-derivative-along-curve} identifies both
derivatives with $D_tZ(t_0)$.
\end{proof}
\section{Geodesics}

Let $M$ be a smooth manifold, possibly with boundary, with a connection $\nabla$
in $TM$. The acceleration of a smooth curve $\gamma\colon I\to M$ is
$D_t\gamma'$. A curve is a geodesic when this field vanishes. In coordinates
$\gamma(t)=(x^1(t),\ldots,x^n(t))$, this is equivalent to
\begin{equation}
\ddot x^k(t)+\dot x^i(t)\dot x^j(t)\Gamma^k_{ij}(x(t))=0. \tag{4.16}
\end{equation}

\begin{definition}[Geodesic]
\label{def:geodesic}
\lean{LeeLib.Ch04.chartAcceleration, LeeLib.Ch04.IsGeodesicInChart, LeeLib.Ch04.isGeodesicInChart_iff}
\uses{thm:covariant-derivative-along-curve}
\leanok



The \emph{acceleration} of $\gamma$ is $D_t\gamma'$. The curve is a
\emph{geodesic} when $D_t\gamma'\equiv0$, equivalently when its velocity is
parallel. Its coordinate functions satisfy (4.16).
\end{definition}

\begin{theorem}[Existence and Uniqueness of Geodesics]
\label{thm:geodesic-existence-uniqueness}
\lean{LeeLib.Ch04.IsSprayGeodesicAt, LeeLib.Ch04.exists_sprayGeodesic_initial, LeeLib.Ch04.sprayGeodesic_eventuallyEq, LeeLib.Ch04.isSprayGeodesicAt_eventuallyEq_of_lift_eq, LeeLib.Ch04.IsSprayGeodesicOnWithInitial.chartAcceleration_zero, LeeLib.Ch04.IsSprayGeodesicOnWithInitial.hasDerivAt_extChartAt_zero}
\dcref{ch4:4.27}

Let $M$ be a smooth manifold without boundary and $\nabla$ a connection in
$TM$. For every $p\in M$, $w\in T_pM$, and $t_0\in\mathbb R$, there are an open
interval $I\ni t_0$ and a geodesic $\gamma\colon I\to M$ such that
\[
\gamma(t_0)=p,
\qquad \gamma'(t_0)=w.
\]
Any two geodesics with these initial data agree on their common domain.
\end{theorem}

\begin{proof}
\sketch In coordinates near $p$, introduce $v^i=\dot x^i$ and rewrite (4.16) as
\begin{align}
\dot x^k(t)&=v^k(t),\\
\dot v^k(t)&=-v^i(t)v^j(t)\Gamma^k_{ij}(x(t)). \tag{4.17}
\end{align}
This is the integral-curve equation for
\begin{equation}
G_{(x,v)}=v^k\frac{\partial}{\partial x^k}\bigg|_{(x,v)}
-v^iv^j\Gamma^k_{ij}(x)\frac{\partial}{\partial v^k}\bigg|_{(x,v)}. \tag{4.18}
\end{equation}
Local flow existence and uniqueness for $G$ give the required geodesic and local
uniqueness. For two solutions, the set of times at which their positions and
velocities agree is nonempty and locally open by the same coordinate argument;
continuity makes it closed in the common interval. Connectedness of the interval
gives equality on the full common domain.
\end{proof}

A geodesic is \emph{maximal} if it has no extension to a strictly larger interval.
A \emph{geodesic segment} has compact interval domain.

\begin{corollary}[Existence of maximal geodesics]
\label{cor:existence-of-maximal-geodesic}
\lean{LeeLib.Ch04.maximalSprayGeodesicInterval, LeeLib.Ch04.maximalSprayGeodesic, LeeLib.Ch04.maximalSprayGeodesicInterval_isOpen, LeeLib.Ch04.maximalSprayGeodesic_structure_of_footInSource}
\uses{thm:geodesic-existence-uniqueness}
\dcref{ch4:4.28}

Let $M$ be a smooth manifold and $\nabla$ a connection in $TM$. For every
$p\in M$ and $v\in T_pM$, there is a unique maximal geodesic
$\gamma\colon I\to M$ on an open interval $I\ni0$ satisfying
$\gamma(0)=p$ and $\gamma'(0)=v$.

\end{corollary}

\begin{proof}
\sketch Take the union of all open intervals containing $0$ on which a geodesic
with the prescribed data exists. Theorem \ref{thm:geodesic-existence-uniqueness}
makes the local geodesics agree on overlaps, so they glue to the unique maximal
one.
\end{proof}

Write $\gamma_v$ for the maximal geodesic with initial velocity $v$ and initial
point $\pi(v)$, where $\pi\colon TM\to M$ is the bundle projection.

\section{Parallel Transport}

A smooth vector or tensor field $V$ along $\gamma$ is \emph{parallel} when
$D_tV=0$. Thus $\gamma$ is geodesic exactly when $\gamma'$ is parallel.

In a coordinate representation $\gamma(t)=(\gamma^1(t),\ldots,\gamma^n(t))$,
a vector field is parallel precisely when
\begin{equation}
\dot V^k(t)=-V^j(t)\dot\gamma^i(t)\Gamma^k_{ij}(\gamma(t)),
\qquad k=1,\ldots,n. \tag{4.19}
\end{equation}
Tensor fields satisfy the analogous linear systems from Proposition 4.18.

\begin{theorem}[Existence, Uniqueness, and Smoothness for Linear ODEs]
\label{thm:linear-ode-existence-uniqueness}
\lean{LeeLib.Ch04.exists_hasDerivWithinAt_Icc, LeeLib.Ch04.IsSolOn, LeeLib.Ch04.flowMap, LeeLib.Ch04.flowMap_injective}
\leanok
\dcref{ch4:4.31}

Let $I\subseteq\mathbb R$ be an open interval and let
$A_i^k\colon I\to\mathbb R$ be smooth for $1\leq i,k\leq n$. For every
$t_0\in I$ and $c=(c^1,\ldots,c^n)\in\mathbb R^n$, the problem
\begin{equation}
\dot V^k(t)=A_i^k(t)V^i(t) \tag{4.20}
\end{equation}
with $V^k(t_0)=c^k$ has a unique smooth solution on all of $I$. The solution
depends smoothly on $(t,c)\in I\times\mathbb R^n$.
\end{theorem}

\begin{proof}
\sketch After translating the parameter, take $t_0=0$. On
$I\times\mathbb R^n$ define
\[
Y=\frac{\partial}{\partial x^0}
+\sum_{k=1}^n A_j^k(x^0)x^j\frac{\partial}{\partial x^k}.
\]
Solutions correspond to integral curves $\eta(t)=(t,V(t))$ with
\begin{equation}
\eta(0)=(0,c^1,\ldots,c^n). \tag{4.21}
\end{equation}
The flow theorem gives local existence, uniqueness, and smooth dependence. To
show the maximal solution fills $I=(a,b)$, suppose its right endpoint is
$b_0<b$. On $[0,b_0]$, choose $M$ with $|A(t)|\leq M$ in Frobenius norm. Then
\[
\frac d{dt}|V|^2=2V^TA(t)V\leq2M|V|^2,
\]
so
\[
\frac d{dt}\left(e^{-2Mt}|V(t)|^2\right)\leq0,
\qquad
|V(t)|\leq e^{Mt}|V(0)|.
\]
Thus $(t,V(t))$ remains in the compact set
$[0,b_0]\times\overline B_R(0)$, with $R=e^{Mb_0}|V(0)|$, contradicting the
escape lemma. Apply the same argument to $-Y$ for the left endpoint, and undo
the parameter translation.
\end{proof}

\begin{theorem}[Existence and Uniqueness of Parallel Transport]
\label{thm:existence-uniqueness-parallel-transport}
\lean{LeeLib.Ch04.exists_isParallelAlongChart_Icc, LeeLib.Ch04.isParallelAlongChart_eqOn_Icc, LeeLib.Ch04.chartParallelTransport}
\uses{thm:covariant-derivative-along-curve, thm:linear-ode-existence-uniqueness}
\leanok
\dcref{ch4:4.32}


Let $M$ be a smooth manifold, possibly with boundary, with a connection in $TM$.
Given a smooth curve $\gamma\colon I\to M$, $t_0\in I$, and either
$v\in T_{\gamma(t_0)}M$ or $v\in T^{(k,l)}_{\gamma(t_0)}M$, there is a unique
parallel vector or tensor field $V$ along $\gamma$ satisfying $V(t_0)=v$.
\end{theorem}

\begin{proof}
\sketch If $\gamma(I)$ lies in one chart, solve (4.19), or the analogous tensor
system, on all of $I$ using Theorem \ref{thm:linear-ode-existence-uniqueness}.
In general, extend the solution through successive chart-covered subintervals.
If the supremum of reachable parameters were an interior point of $I$, a chart
around its image and the single-chart result would extend the field farther.
The same argument to the left of $t_0$ gives existence on all of $I$; local
uniqueness propagates across overlaps.
\end{proof}

The resulting field is the \emph{parallel transport} of $v$ along $\gamma$. For
$t_0,t_1\in I$, define
\begin{equation}
P_{t_0t_1}^\gamma\colon T_{\gamma(t_0)}M\longrightarrow T_{\gamma(t_1)}M,
\qquad P_{t_0t_1}^\gamma(v)=V(t_1). \tag{4.22}
\end{equation}
This map is linear, and $P_{t_1t_0}^\gamma$ is its inverse.

For an admissible curve $\gamma\colon[a,b]\to M$, a \emph{piecewise smooth vector
field along $\gamma$} is a continuous $V\colon[a,b]\to TM$, with
$V(t)\in T_{\gamma(t)}M$, that is smooth on every subinterval of some admissible
partition. It is parallel when $D_tV=0$ wherever both fields are smooth.

\begin{corollary}[Parallel Transport Along Piecewise Smooth Curves]
\label{cor:parallel-transport-piecewise-smooth-curves}
\lean{LeeLib.Ch04.exists_isPiecewiseParallelAlongChart_Icc, LeeLib.Ch04.isPiecewiseParallelAlongChart_eqOn_Icc}
\leanok
\dcref{ch4:4.33}

Let $M$ be a smooth manifold, possibly with boundary, with a connection in $TM$.
For an admissible curve $\gamma\colon[a,b]\to M$ and initial vector
$v\in T_{\gamma(a)}M$ or tensor $v\in T^{(k,l)}(T_{\gamma(a)}M)$, there is a
unique piecewise smooth parallel field $V$ along $\gamma$ with $V(a)=v$.
Moreover, $V$ is smooth wherever $\gamma$ is smooth.
\end{corollary}

\begin{proof}
\sketch On an admissible partition $(a_0,\ldots,a_k)$, transport $v$ along the
first smooth segment and use its endpoint value as the initial value on the next
segment. Induction gives existence, and segmentwise uniqueness gives global
uniqueness.
\end{proof}

Transporting a basis of $T_{\gamma(t_0)}M$ produces a \emph{parallel frame}
$(E_1,\ldots,E_n)$ along $\gamma$. Since parallel transport is an isomorphism,
this is a basis at each time. If $V=V^iE_i$, then
\begin{equation}
D_tV(t)=\dot V^i(t)E_i(t) \tag{4.23}
\end{equation}
wherever $V$ and $\gamma$ are smooth. Hence $V$ is parallel exactly when its
components in a parallel frame are constant.

\begin{theorem}[Parallel Transport Determines Covariant Differentiation]
\label{thm:parallel-transport-determines-covariant-diff}
\uses{cor:parallel-transport-piecewise-smooth-curves}
\dcref{ch4:4.34}

For a smooth curve $\gamma\colon I\to M$, a smooth vector field $V$ along it,
and $t_0\in I$,
\begin{equation}
D_tV(t_0)=\lim_{t_1\to t_0}
\frac{P_{t_1t_0}^\gamma V(t_1)-V(t_0)}{t_1-t_0}. \tag{4.24}
\end{equation}
\end{theorem}

\begin{proof}
\sketch Write $V(t)=V^i(t)E_i(t)$ in a parallel frame. Equation (4.23) gives
$D_tV(t_0)=\dot V^i(t_0)E_i(t_0)$, while
\[
P_{t_1t_0}^\gamma V(t_1)=V^i(t_1)E_i(t_0).
\]
The difference quotient has the same limit.
\end{proof}

\begin{corollary}[Parallel Transport Determines the Connection]
\label{cor:parallel-transport-determines-connection}
\uses{thm:parallel-transport-determines-covariant-diff}
\dcref{ch4:4.35}

Let $X,Y$ be smooth vector fields. For $p\in M$ and any smooth curve
$\gamma\colon I\to M$ satisfying $\gamma(0)=p$ and $\gamma'(0)=X_p$,
\begin{equation}
\nabla_XY|_p=\lim_{h\to0}
\frac{P_{h0}^\gamma Y(\gamma(h))-Y_p}{h}. \tag{4.25}
\end{equation}
\end{corollary}

\begin{proof}
\sketch Restrict $Y$ to $\gamma$. The defining property of covariant
differentiation along curves gives $\nabla_XY|_p=D_t(Y\circ\gamma)(0)$; apply
Theorem \ref{thm:parallel-transport-determines-covariant-diff}.
\end{proof}

A smooth vector or tensor field on $M$ is \emph{parallel} when it is parallel
along every smooth curve.

\begin{proposition}
\label{prop:parallel-field-characterization}
\dcref{ch4:4.36}
Let $M$ be a smooth manifold, possibly with boundary, with a connection $\nabla$
in $TM$. A smooth vector or tensor field $A$ is parallel on $M$ if and only if
$\nabla A\equiv0$.
\end{proposition}

\section{Pullback Connections}

For a diffeomorphism $\varphi\colon M\to\widetilde M$, write $\varphi_*X$ for
the vector field satisfying
$d\varphi_p(X_p)=(\varphi_*X)_{\varphi(p)}$.

\begin{lemma}[Pullback Connections]
\label{lem:pullback-connection}
\lean{LeeLib.Ch04.TangentConnection.pullback}
\uses{lem:pushforward-vector-field}
\leanok
\dcref{ch4:4.37}



Let $M,\widetilde M$ be smooth manifolds, possibly with boundary, let
$\widetilde\nabla$ be a connection in $T\widetilde M$, and let
$\varphi\colon M\to\widetilde M$ be a diffeomorphism. Then
\begin{equation}
(\varphi^*\widetilde\nabla)_XY
=(\varphi^{-1})_*
\left(\widetilde\nabla_{\varphi_*X}(\varphi_*Y)\right) \tag{4.26}
\end{equation}
defines a connection in $TM$, called the \emph{pullback connection}.
\end{lemma}

\begin{proof}
\sketch Real linearity in $Y$ is immediate. For $f\in C^\infty(M)$, set
$\widetilde f=f\circ\varphi^{-1}$. Since
$\varphi_*(fX)=\widetilde f\,\varphi_*X$, equation (4.26) gives
\[
(\varphi^*\widetilde\nabla)_{fX}Y
=f(\varphi^*\widetilde\nabla)_XY.
\]
Also $(\varphi_*X)(\widetilde f)=(Xf)\circ\varphi^{-1}$, so the Leibniz rule for
$\widetilde\nabla$ yields
\[
(\varphi^*\widetilde\nabla)_X(fY)
=f(\varphi^*\widetilde\nabla)_XY+(Xf)Y.
\]
Thus (4.26) satisfies the connection axioms.
\end{proof}

\begin{proposition}[Properties of Pullback Connections]
\label{prop:pullback-connection-properties}
\lean{LeeLib.Ch04.covariantDeriv_pullback_naturality, LeeLib.Ch04.pullback_covariantDeriv_apply}
\uses{lem:pullback-connection}
\dcref{ch4:4.38}


Let $\nabla=\varphi^*\widetilde\nabla$ and let $\gamma\colon I\to M$ be smooth.
\begin{enumerate}
\item[(a)] For every smooth vector field $V$ along $\gamma$,
\[
d\varphi\circ D_tV=\widetilde D_t(d\varphi\circ V),
\]
where the two derivatives are taken along $\gamma$ and $\varphi\circ\gamma$.
\item[(b)] If $\gamma$ is a $\nabla$-geodesic, then $\varphi\circ\gamma$ is a
$\widetilde\nabla$-geodesic.
\item[(c)] For $t_0,t_1\in I$,
\[
d\varphi_{\gamma(t_1)}\circ P_{t_0t_1}^\gamma
=P_{t_0t_1}^{\varphi\circ\gamma}\circ d\varphi_{\gamma(t_0)}.
\]
\end{enumerate}
\end{proposition}
\section{Further Results}

\begin{proposition}
\label{prop:embedded-submanifold-covariant-derivative}
\dcref{ch4:4-1}
\uses{prop:tangent-space-submanifold}
Let $M\subseteq\mathbb R^n$ be an embedded submanifold, let
$Y\in\mathfrak X(M)$, and define $\nabla_v^\top Y$ by (4.4) for
$p\in M$ and $v\in T_pM$.
\begin{enumerate}
\item[(a)] $\nabla_v^\top Y$ is independent of the chosen extension
$\widetilde Y$ of $Y$.
\item[(b)] The construction is invariant under rigid motions: if $F\in E(n)$ and
$\widetilde M=F(M)$, then
\[
dF_p(\nabla_v^\top Y)=\nabla_{dF_p(v)}^\top(F_*Y).
\]
\end{enumerate}
\end{proposition}

\begin{exercise}
\label{prob:embedded-submanifold-covariant-derivative}
\dcref{ch4:4-1}
\uses{prop:tangent-space-submanifold}
Let $M \subseteq \mathbb{R}^n$ be an embedded submanifold and $Y \in \mathfrak{X}(M)$. For every point $p \in M$ and vector $v \in T_p M$, define $\nabla_v^\top Y$ by (4.4).

(a) Show that $\nabla_v^\top Y$ does not depend on the choice of extension $\tilde{Y}$ of $Y$. [Hint: Use Prop. A.28.]

(b) Show that $\nabla_v^\top Y$ is invariant under rigid motions of $\mathbb{R}^n$, in the following sense: if $F \in E(n)$ and $\tilde{M} = F(M)$, then $dF_p(\nabla_v^\top Y) = \nabla_{dF_p(v)}^\top (F_* Y)$.

(Used on pp. 87, 93.)
\end{exercise}


\begin{proposition}
\label{prop:torsion-tensor}
\dcref{ch4:4-6}
\uses{ex:covariant-hessian}

Let $\nabla$ be a connection in $TM$ and define
\[
\tau(X,Y)=\nabla_XY-\nabla_YX-[X,Y].
\]
\begin{enumerate}
\item[(a)] $\tau$ is a $(1,2)$-tensor field, called the
\emph{torsion tensor} of $\nabla$.
\item[(b)] Calling $\nabla$ \emph{symmetric} when $\tau\equiv0$, this is
equivalent to
\[
\Gamma^k_{ij}=\Gamma^k_{ji}
\]
in every coordinate frame. \emph{Warning.} The coefficients need not be symmetric
in a noncoordinate frame.
\item[(c)] The connection $\nabla$ is symmetric if and only if the covariant Hessian
$\nabla^2u$ is a symmetric $2$-tensor for every $u\in C^\infty(M)$.
\item[(d)] The Euclidean connection on $\mathbb R^n$ is symmetric.
\end{enumerate}
\end{proposition}

\begin{proposition}[Two connections and their difference tensor]
\label{prop:connections-difference-tensor}
\dcref{ch4:4-9}
\uses{prop:torsion-tensor}

Let $\nabla^0,\nabla^1$ be connections on $TM$, with difference tensor $D$.
\begin{enumerate}
\item[(a)] The two connections have the same torsion if and only if
$D$ is symmetric:
\[
D(X,Y)=D(Y,X).
\]
\item[(b)] They determine the same geodesics if and only if $D$ is
antisymmetric:
\[
D(X,Y)=-D(Y,X).
\]
\end{enumerate}
\end{proposition}

\begin{proposition}[Parallel vector fields and reparametrization]
\label{prop:parallel-reparametrization}
\dcref{ch4:4-10}

Let $M$ carry a connection, let $\gamma\colon I\to M$ be smooth, and let
$Y\in\mathfrak X(\gamma)$. If $Y$ is parallel along $\gamma$, then its induced
field is parallel along every reparametrization of $\gamma$.
\end{proposition}

\begin{exercise}
\label{exer:complete-proof-lemma-4.10}
\dcref{ch4:4.11}
\uses{lem:connections-from-connection-coefficients}
% The step Lee leaves as an exercise — that the formula (4.9) definition satisfies the
% connection axioms — is completed: connOfCoeff_isCovariantDerivativeOn proves the
% additivity (`add`) and the product rule (`leibniz`) directly from (4.9) on the global
% trivialization frame, and connectionCoeff_connOfCoeff recovers the coefficients.
\lean{LeeLib.Ch04.connOfCoeff_isCovariantDerivativeOn, LeeLib.Ch04.connectionCoeff_connOfCoeff}
\leanok
Complete the proof of Lemma 4.10.
\end{exercise}

\begin{exercise}\label{exer:connection-restriction-is-connection}\dcref{ch4:4.4}
\uses{prop:connection-restriction}
Complete the proof of the preceding proposition by showing that $\nabla^U$ is a connection.
\end{exercise}

\begin{exercise}
\label{exer:covariant-derivative-properties}
\dcref{ch4:4.25}
\uses{thm:covariant-derivative-along-curve}
Complete the proof of Theorem 4.24 by showing that the operator $D_t$ defined in coordinates by (4.15) satisfies properties (i)--(iii).
\end{exercise}

\begin{exercise}[Geodesics in Euclidean space]
\label{exer:euclidean-geodesics}
\dcref{ch4:4.29}
\uses{prop:connection-restriction}
Show that the maximal geodesics on $\mathbb{R}^n$ with respect to the Euclidean connection (4.3) are exactly the constant curves and the straight lines with constant-speed parametrizations.
\end{exercise}

\begin{exercise}\label{exer:lemma-4-1-comparison}\dcref{ch4:4.2}
\uses{lem:connection-is-local-operator}
Complete the proof of Lemma 4.1 by showing that $\nabla_X Y$ and $\nabla_{\bar{X}} Y$ agree at $p$ if $X = \bar{X}$ on a neighborhood of $p$.
\end{exercise}

\begin{exercise}\label{exer:parallel-vector-field-euclidean-connection}\dcref{ch4:4.30}
Let $\gamma : I \to \mathbb{R}^n$ be a smooth curve, and let $V$ be a smooth vector field along $\gamma$. Show that $V$ is parallel along $\gamma$ with respect to the Euclidean connection if and only if its component functions (with respect to the standard basis) are constants.
\end{exercise}

\begin{exercise}
\label{exer:prove-proposition-4-18}
\dcref{ch4:4.19}
\uses{prop:total-covariant-derivative-components}
Prove Proposition \ref{prop:total-covariant-derivative-components}.
\end{exercise}

\begin{exercise}
\label{exer:tensor-product-covariant-derivative}
\dcref{ch4:4.20}
Suppose $F$ is a smooth $(k,l)$-tensor field and $G$ is a smooth $(r,s)$-tensor field. Show that the components of the total covariant derivative of $F \otimes G$ are given by
\[
(\nabla(F \otimes G))^{j_1\cdots j_k p_1\cdots p_r}_{i_1\cdots i_l q_1\cdots q_s;m} = F^{j_1\cdots j_k}_{i_1\cdots i_l;m} G^{p_1\cdots p_r}_{q_1\cdots q_s} + F^{j_1\cdots j_k}_{i_1\cdots i_l} G^{p_1\cdots p_r}_{q_1\cdots q_s;m}.
\]

[Remark: This formula is often written in the following way, more suggestive of the product rule for ordinary derivatives:
\[
\left(F^{j_1\cdots j_k}_{i_1\cdots i_l} G^{p_1\cdots p_r}_{q_1\cdots q_s}\right)_{;m} = F^{j_1\cdots j_k}_{i_1\cdots i_l;m} G^{p_1\cdots p_r}_{q_1\cdots q_s} + F^{j_1\cdots j_k}_{i_1\cdots i_l} G^{p_1\cdots p_r}_{q_1\cdots q_s;m}.
\]

Notice that this does not say that $\nabla(F \otimes G) = (\nabla F) \otimes G + F \otimes (\nabla G)$, because in the first term on the right-hand side of this latter formula, the index resulting from differentiation is not the last lower index.]
\end{exercise}

\begin{exercise}\label{prob:characterizing-connections}\dcref{ch4:4-4}
\uses{thm:connections-affine-space}
Prove Theorem 4.14 (characterizing the space of connections).
\end{exercise}

\begin{exercise}[Connection 1-forms and Cartan's first structure equation]
\label{prob:connection-1forms-cartan}
\dcref{ch4:4-14}
\uses{prob:torsion-tensor}
Let $M$ be a smooth $n$-manifold and $\nabla$ a connection in $TM$, let $(E_i)$ be a local frame on some open subset $U \subset M$, and let $(\epsilon^i)$ be the dual coframe.

(a) Show that there is a uniquely determined $n \times n$ matrix of smooth 1-forms $(\omega^i_j)$ on $U$, called the \textbf{connection 1-forms} for this frame, such that
\[
\nabla_X E_i = \omega^i_j(X) E_j
\]
for all $X \in \mathfrak{X}(U)$.

(b) \textbf{Cartan's first structure equation}: Prove that these forms satisfy the following equation, due to \'Elie Cartan:
\[
d\epsilon^i = \epsilon^i \wedge \omega^i_j + \tau^i,
\]
where $\tau^1, \ldots, \tau^n \in \Omega^2(M)$ are the \textbf{torsion 2-forms}, defined in terms of the torsion tensor $\tau$ (Problem 4-6) and the frame $(E_i)$ by
\[
\tau(X,Y) = \tau^i(X,Y) E_i.
\]

(Used on pp. 145, 222.)
\end{exercise}

\begin{exercise}[Two connections and their difference tensor]
\label{prob:connections-difference-tensor}
\dcref{ch4:4-9}
\uses{prob:torsion-tensor}
% Part (a) formalized (sameTorsion_iff_differenceTensor_symm), via the identity
% τ¹(X,Y)−τ⁰(X,Y) = D(Y,X)−D(X,Y).
% Part (b) formalized at the level of the geodesic equation
% (sameChartAcceleration_iff_chartGamma_diff_antisymm,
% isGeodesicInChart_congr_of_diff_antisymm): the two connections induce the same
% acceleration D_tγ' on every chart curve — hence the same chart geodesics — iff the
% difference Γ¹−Γ⁰ of their chart Christoffel contractions is antisymmetric, via the
% polarization identity biadditive_diag_eq_zero_iff_antisymm.  This chart condition is
% bridged to Lee's abstract difference tensor D = ∇¹−∇⁰ (differenceTensor, Prop 4.13):
% chartGamma_diff_antisymm_iff_differenceTensor_frame_antisymm shows Γ¹−Γ⁰ is
% antisymmetric iff D is antisymmetric on the chart frame, D(∂_j,∂_i)=−D(∂_i,∂_j) — the
% exact geodesic counterpart of part (a) (torsion ↔ D symmetric), via the frame identity
% differenceTensor_chartFrame_eq_sum (D's chart components = Christoffel coefficient
% differences).  The equivalence with the same solution SET (rather than the same
% equation) is the spray existence/uniqueness theory, still partial in the chapter, so
% no \leanok on the whole problem.
\lean{LeeLib.Ch04.sameTorsion_iff_differenceTensor_symm, LeeLib.Ch04.sameChartAcceleration_iff_chartGamma_diff_antisymm, LeeLib.Ch04.isGeodesicInChart_congr_of_diff_antisymm, LeeLib.Ch04.chartGamma_diff_antisymm_iff_differenceTensor_frame_antisymm, LeeLib.Ch04.differenceTensor_chartFrame_eq_sum}
Let $M$ be a smooth manifold, and let $\nabla^0$ and $\nabla^1$ be two connections on $TM$.

(a) Show that $\nabla^0$ and $\nabla^1$ have the same torsion (Problem 4-6) if and only if their difference tensor is symmetric, i.e., $D(X,Y) = D(Y,X)$ for all $X$ and $Y$.

(b) Show that $\nabla^0$ and $\nabla^1$ determine the same geodesics if and only if their difference tensor is antisymmetric, i.e., $D(X,Y) = -D(Y,X)$ for all $X$ and $Y$.

(Used on p. 145.)
\end{exercise}

\begin{exercise}\label{prob:figure-eight-immersion}\dcref{ch4:4-7}
\uses{ex:figure-eight-curve}
Let $\gamma: (-\pi, \pi) \to \mathbb{R}^2$ be the figure eight curve defined in Example 4.23. Prove that $\gamma$ is an injective smooth immersion, but its velocity vector field is not extendible.
\end{exercise}

\begin{exercise}\label{prob:lie-derivative-not-connection}\dcref{ch4:4-2}
\uses{prop:lie-derivative-equals-bracket}
In your study of smooth manifolds, you have already seen another way of taking ``directional derivatives of vector fields,'' the Lie derivative $\mathcal{L}_X Y$ (which is equal to the Lie bracket $[X,Y]$; see Prop.~A.46). Suppose $M$ is a smooth manifold of positive dimension.

(a) Show that the map $\mathcal{L}: \mathfrak{X}(M) \times \mathfrak{X}(M) \to \mathfrak{X}(M)$ is not a connection.

(b) Show that there are smooth vector fields $X$ and $Y$ on $\mathbb{R}^2$ such that $X = Y = \partial_1$ along the $x^1$-axis, but the Lie derivatives $\mathcal{L}_X(\partial_2)$ and $\mathcal{L}_Y(\partial_2)$ are not equal on the $x^1$-axis.
\end{exercise}

\begin{exercise}[Connection on Lie groups]
\label{prob:lie-group-connection}
\dcref{ch4:4-11}
\uses{prob:torsion-tensor}
Suppose $G$ is a Lie group.

(a) Show that there is a unique connection $\nabla$ in $TG$ with the property that every left-invariant vector field is parallel.

(b) Show that the torsion tensor of $\nabla$ (Problem 4-6) is zero if and only if $G$ is abelian.
\end{exercise}

\begin{exercise}[Parallel vector fields and reparametrization]
\label{prob:parallel-reparametrization}
\dcref{ch4:4-10}
% Formalized in chart coordinates (isParallelAlongChart_comp), via the chain-rule
% rescaling of the covariant derivative along a curve under reparametrization
% D_s(V∘φ) = φ' • (D_t V)∘φ (covariantDerivAlongChart_comp); parallel (D_t V ≡ 0)
% is therefore preserved.  Chart-local like the rest of the D_t development, so no
% \leanok on the abstract 𝔛(γ) statement.
\lean{LeeLib.Ch04.covariantDerivAlongChart_comp, LeeLib.Ch04.isParallelAlongChart_comp}
Suppose $M$ is a smooth manifold endowed with a connection, $\gamma: I \to M$ is a smooth curve, and $Y \in \mathfrak{X}(\gamma)$. Prove that if $Y$ is parallel along $\gamma$, then it is parallel along every reparametrization of $\gamma$.
\end{exercise}

\begin{exercise}[Parallel tensor fields]
\label{prob:parallel-tensor-fields}
\dcref{ch4:4-12}
\uses{prop:parallel-field-characterization}
Prove Proposition 4.36 (a vector or tensor field $A$ is parallel if and only if $\nabla A = 0$).
\end{exercise}

\begin{exercise}[Pullback connections]
\label{prob:pullback-connections}
\dcref{ch4:4-13}
\uses{prop:pullback-connection-properties}
Prove Proposition 4.38 (properties of pullback connections).
\end{exercise}

\begin{exercise}\label{prob:tensor-field-formulas}\dcref{ch4:4-5}
\uses{prop:covariant-derivative-components}
Prove Proposition 4.16 (local formulas for covariant derivatives of tensor fields).
\end{exercise}

\begin{exercise}\label{prob:torsion-tensor}\dcref{ch4:4-6}
\uses{ex:covariant-hessian}
\lean{LeeLib.Ch04.torsion, LeeLib.Ch04.torsion_apply, LeeLib.Ch04.isSymmetric_iff}
Let $M$ be a smooth manifold and let $\nabla$ be a connection in $TM$. Define a map $\tau: \mathfrak{X}(M) \times \mathfrak{X}(M) \to \mathfrak{X}(M)$ by
\[
\tau(X,Y) = \nabla_X Y - \nabla_Y X - [X,Y].
\]

(a) Show that $\tau$ is a $(1,2)$-tensor field, called the \textit{torsion tensor} of $\nabla$.

(b) We say that $\nabla$ is \textit{symmetric} if its torsion vanishes identically. Show that $\nabla$ is symmetric if and only if its connection coefficients with respect to every coordinate frame are symmetric: $\Gamma^k_{ij} = \Gamma^k_{ji}$. [Warning: They might not be symmetric with respect to other frames.]

(c) Show that $\nabla$ is symmetric if and only if the covariant Hessian $\nabla^2 u$ of every smooth function $u \in C^\infty(M)$ is a symmetric 2-tensor field. (See Example 4.22.)

(d) Show that the Euclidean connection $\nabla$ on $\mathbb{R}^n$ is symmetric.

(Used on pp. 113, 121, 123.)
\end{exercise}

\begin{exercise}\label{prob:transformation-law-connection}\dcref{ch4:4-3}
\uses{prop:transformation-law-connection-coefficients}
\lean{LeeLib.Ch04.connectionCoeff_transformation_law}
\leanok
Prove Proposition 4.7 (the transformation law for the connection coefficients).
\end{exercise}

\begin{exercise}\label{prob:vector-field-extendibility}\dcref{ch4:4-8}
Suppose $M$ is a smooth manifold (without boundary), $I \subseteq \mathbb{R}$ is an interval (bounded or not, with or without endpoints), and $\gamma: I \to M$ is a smooth curve.

(a) Show that for every $t_0 \in I$ such that $\gamma'(t_0) \neq 0$, there is a connected neighborhood $J$ of $t_0$ in $I$ such that every smooth vector field along $\gamma|_J$ is extendible.
\end{exercise}
